\subsection{Expanding molecular orbitals}

In practical molecular calculations, each molecular orbital is expanded in a
finite set of known one-electron functions,
\begin{equation}
  \boxed{\psi_p(\br)=\sum_{\mu=1}^{K}C_{\mu p}\phi_\mu(\br)}.
\end{equation}
The basis functions $\phi_\mu$ are fixed, whereas the coefficients $C_{\mu p}$
are determined variationally. Increasing $K$ systematically enlarges the
variational space and moves the result toward the complete-basis-set (CBS)
limit. Consequently, the basis set is the only source of empirical parameters
in an otherwise ab initio calculation.

A useful basis must be accurate, compact, numerically stable, and practical.
These requirements compete: a larger and more flexible basis improves the
description of the orbitals, but increases both the computational cost and the
risk of linear dependence.

\paragraph{The LEGO principle.}

Basis-set construction follows a LEGO-like philosophy: a small collection of
well-designed building blocks is combined to describe chemically diverse
systems.
\begin{figure}[htbp]
  \centering
  \includegraphics[width=0.62\textwidth]{lego}
  \caption{The basis functions are the elementary building blocks from which
  molecular orbitals are assembled.}
\end{figure}

\subsubsection{Common basis functions}

Several families of one-electron functions are routinely used.
\begin{itemize}
  \item \red{Gaussian-type orbitals (GTOs)}, proportional to
  $\exp(-\alpha r^2)$, are atom centered and are the standard choice for
  molecules. Their multicenter integrals are efficient. They are used in, for
  example, \textsc{Orca}, \textsc{Q-Chem}, \textsc{Molcas}, \textsc{Molpro},
  \textsc{PySCF}, \textsc{Psi4}, \textsc{CFOUR}, \textsc{MRCC}, and
  \textsc{GAMESS}.
  \item \blue{Slater-type orbitals (STOs)}, proportional to
  $\exp(-\zeta r)$, are atom centered and have the correct nuclear cusp and
  long-range tail. They are used, for example, in \textsc{ADF}.
  \item \orange{Plane waves (PWs)}, proportional to $\exp(\ii\bk\cdot\br)$,
  are delocalized, mutually orthonormal, and naturally adapted to periodic
  boundary conditions. Their size is controlled by an energy cutoff. They are
  used in \textsc{Quantum Espresso}, \textsc{VASP}, \textsc{ABINIT}, and
  \textsc{CASTEP}.
  \item \green{Numerical atomic orbitals (NAOs)} have tabulated radial parts.
  They are compact and flexible and are used in \textsc{FHI-aims} and
  \textsc{SIESTA}.
\end{itemize}

\subsection{Plane waves}

In a unit cell of volume $\Omega$, a normalized plane wave is
\begin{equation}
  \phi_{\bk}(\br)=\frac{1}{\sqrt{\Omega}}\exp(\ii\bk\cdot\br),
  \qquad \abs{\bk}\leq k_\text{cut},
  \qquad E_\text{cut}=\frac{k_\text{cut}^2}{2}.
\end{equation}
The reciprocal-lattice vector is
\begin{equation}
  \bk=k_1\bb_1+k_2\bb_2+k_3\bb_3,
\end{equation}
where, for direct-lattice vectors $\ba_1,\ba_2,\ba_3$,
\begin{align}
  \Omega&=\ba_1\cdot(\ba_2\times\ba_3),
  &\bA&=(\ba_1,\ba_2,\ba_3),\\
  \bb_1&=2\pi\frac{\ba_2\times\ba_3}{\Omega},
  &\bb_2&=2\pi\frac{\ba_3\times\ba_1}{\Omega},
  &\bb_3&=2\pi\frac{\ba_1\times\ba_2}{\Omega}.
\end{align}

The overlap integrals are diagonal:
\begin{align}
  \braket{\bk_1}{\bk_2}
  &=\frac{1}{\Omega}\int_\Omega
    \exp[-\ii\bk_1\cdot\br]\exp[\ii\bk_2\cdot\br]\dd\br\\
  &=\frac{1}{\Omega}\int_\Omega
    \exp[\ii(\bk_2-\bk_1)\cdot\br]\dd\br
  =\delta_{\bk_1\bk_2}.
\end{align}
The kinetic-energy integrals are diagonal as well:
\begin{align}
  \mel{\bk_1}{-\frac12\nabla^2}{\bk_2}
  &=\frac{1}{\Omega}\int_\Omega
    \exp[-\ii\bk_1\cdot\br]
    \qty(-\frac12\nabla^2)\exp[\ii\bk_2\cdot\br]\dd\br\\
  &=\frac{k_2^2}{2\Omega}\int_\Omega
    \exp[\ii(\bk_2-\bk_1)\cdot\br]\dd\br
  =\frac{k_2^2}{2}\delta_{\bk_1\bk_2}.
\end{align}
The electron-repulsion integral begins as
\begin{align}
  \rbraket{\bk_1\bk_2}{\bk_3\bk_4}
  &=\frac{1}{\Omega^2}\iint_\Omega
  \frac{\exp[\ii(\bk_3-\bk_1)\cdot\br_1]
        \exp[\ii(\bk_4-\bk_2)\cdot\br_2]}
       {r_{12}}\dd\br_1\dd\br_2.
\end{align}
Using the reciprocal-space resolution of the Coulomb operator,
\begin{equation}
  \frac{1}{r_{12}}=\frac{1}{\Omega}\sum_{\bk\ne\bO}
  \frac{4\pi}{k^2}\exp[\ii\bk\cdot(\br_1-\br_2)],
\end{equation}
where the $\bk=\bO$ term is delicate but can be treated, gives
\begin{align}
  \rbraket{\bk_1\bk_2}{\bk_3\bk_4}
  &=\frac{1}{\Omega^3}\sum_{\bk\ne\bO}\frac{4\pi}{k^2}
  \iint_\Omega\exp[\ii(\bk_3-\bk_1-\bk)\cdot\br_1]
  \exp[\ii(\bk_4-\bk_2+\bk)\cdot\br_2]\dd\br_1\dd\br_2\\
  &=\frac{1}{\Omega}\sum_{\bk\ne\bO}\frac{4\pi}{k^2}
  \delta_{\bk_3-\bk_1-\bk,\bO}
  \delta_{\bk_4-\bk_2+\bk,\bO}\\
  &=\frac{4\pi}{\Omega\abs{\bk_1-\bk_3}^2}
  \delta_{\bk_1+\bk_2,\bk_3+\bk_4}.
\end{align}
The Kronecker delta expresses conservation of crystal momentum.

The key physical picture is that assemblies of atoms are treated as slight
distortions of free electrons. Plane waves are orthonormal and independent of
the atomic positions; naturally compatible with periodic boundary conditions;
systematically convergent through the single parameter $E_\text{cut}$; and
diagonalize the kinetic and Coulomb operators in reciprocal space. They also
permit efficient fast-Fourier-transform evaluation and have neither basis-set
superposition error nor atom-centered Pulay forces. Their limitations are
large basis sizes,
\begin{equation}
  N_\text{PW}\propto\Omega E_\text{cut}^{3/2}.
\end{equation}
Core orbitals require prohibitively large cutoffs, motivating
pseudopotentials; the cost increases with supercell size; delocalized functions
provide little real-space sparsity; and the large virtual spaces make
correlated methods expensive.
Plane waves are therefore especially effective for periodic valence-electron
calculations, usually in conjunction with pseudopotentials.

\subsection{Slater and Gaussian orbitals}

The general Slater- and Gaussian-type orbitals used in the comparison are
\begin{equation}
 \boxed{\phi_{\zeta n\ell m}^\text{STO}(\br)
 \propto r^{n-1}e^{-\zeta r}Y_{\ell m}(\hat{\br})},
 \qquad
 \boxed{\phi_{\alpha n\ell m}^\text{GTO}(\br)
 \propto r^{2(n-1)-\ell}e^{-\alpha r^2}Y_{\ell m}(\hat{\br})}.
\end{equation}
For their $1s$ representatives, the STO has the correct nuclear cusp,
\begin{equation}
 \phi_\text{1s}^\text{STO}(r)\propto e^{-\zeta r},
 \qquad
 \eval{\dv{\phi_\text{1s}^\text{STO}(r)}{r}}_{r=0}=-\zeta,
\end{equation}
whereas the GTO is flat at the nucleus,
\begin{equation}
 \phi_\text{1s}^\text{GTO}(r)\propto e^{-\alpha r^2},
 \qquad
 \eval{\dv{\phi_\text{1s}^\text{GTO}(r)}{r}}_{r=0}=0.
\end{equation}
The STO also has the correct exponential long-range behavior of a molecular
orbital,
\begin{equation}
 \phi_\text{HOMO}^\text{STO}(r)
 \propto e^{-\sqrt{-2\epsilon_\text{HOMO}}r},
 \qquad \mathrm{IP}=-\epsilon_\text{HOMO}.
\end{equation}
A Gaussian instead decays too rapidly. Multicenter STO integrals are hard to
compute, whereas multicenter GTO integrals are easy. GTOs usually win, but more
of them are required.

\subsubsection{Primitive and contracted Gaussians}

A primitive Cartesian Gaussian centered at $\bA$ is
\begin{equation}
  \sket{\ba}=(x-A_x)^{a_x}(y-A_y)^{a_y}(z-A_z)^{a_z}
  \exp[-\alpha\abs{\br-\bA}^2].
\end{equation}
A contracted Gaussian is a fixed linear combination of primitives,
\begin{equation}
  \boxed{
  \ket{\ba}\equiv\phi_{\ba}^{\bA}(\br)
  =\sum_{k=1}^{K}D_k\sket{\ba_k}}.
\end{equation}
Here $\alpha$ is the exponent, $\bA$ the center, and
$\ba=(a_x,a_y,a_z)$ the angular-momentum vector, with total angular momentum
$a=a_x+a_y+a_z$.

An $s$ shell contains $(0,0,0)$; a $p$ shell contains $(1,0,0)$,
$(0,1,0)$, and $(0,0,1)$; a Cartesian $d$ shell contains
$(2,0,0)$, $(0,2,0)$, $(0,0,2)$, $(1,1,0)$, $(1,0,1)$, and $(0,1,1)$;
and a Cartesian $f$ shell contains
\begin{equation}
 \begin{gathered}
 (3,0,0),(0,3,0),(0,0,3),(2,1,0),(2,0,1),\\
 (1,2,0),(0,2,1),(0,1,2),(1,0,2),(1,1,1).
 \end{gathered}
\end{equation}
In general,
\begin{equation}
  N_\ell^\text{cart}=\frac{(\ell+1)(\ell+2)}{2},
\end{equation}
which gives $1,3,6,10$ functions for $s,p,d,f$. Spherical harmonics instead
give $N_\ell^\text{sph}=2\ell+1$, namely $1,3,5,7$. Thus
\begin{equation}
  6d=5d+1s,
  \qquad
  10f=7f+3p,
\end{equation}
because, for example, $x^2+y^2+z^2=r^2$. Most molecular calculations use
pure spherical functions.

Contraction reduces the number of variational coefficients, but it does not
remove the primitive-integral work. For four contracted functions,
\begin{equation}
  \rbraket{\ba\bb}{\bc\bd}
  =\sum_{i=1}^{K_a}\sum_{j=1}^{K_b}\sum_{k=1}^{K_c}\sum_{l=1}^{K_d}
  D_i^aD_j^bD_k^cD_l^d
  \sbraket{\ba_i\bb_j}{\bc_k\bd_l}.
\end{equation}
For example, a cc-pVTZ carbon $1s$ function contracted from ten primitives
requires $10^4$ primitive $ssss$ integrals for a single four-center integral.
\begin{figure}[htbp]
  \centering
  \includegraphics[width=0.62\textwidth]{C}
  \caption{Contracted Gaussian functions for carbon.}
\end{figure}
Basis-set definitions can be inspected and downloaded from the
\href{https://www.basissetexchange.org}{Basis Set Exchange}.

\subsection{Choosing a Gaussian basis}

Different physical effects demand different forms of flexibility.
\begin{itemize}
  \item \red{Radial flexibility} is provided by multiple-$\zeta$ functions,
  for example double-$\zeta$, triple-$\zeta$, and quadruple-$\zeta$ bases.
  \item \blue{Angular flexibility} is supplied by polarization functions,
  such as $d$ functions on second-row atoms and $p$ functions on hydrogen.
  \item \orange{Long-range flexibility} is supplied by diffuse functions
  with small exponents, which are important for anions, Rydberg states,
  noncovalent interactions, multipoles, and (hyper)polarizabilities.
  \item \green{Core flexibility} is supplied by tight functions and is needed
  for core correlation, nuclear magnetic shielding, and spin--spin coupling
  constants.
\end{itemize}

\begin{figure}[htbp]
  \centering
  \includegraphics[width=0.72\textwidth]{BSE}
  \caption{The Basis Set Exchange provides standardized definitions for a
  large number of orbital and auxiliary basis sets.}
\end{figure}

\paragraph{STO-$n$G.}

In an STO-$n$G basis, there is one contracted basis function for every occupied
orbital in the minimal atomic description. Each contracted GTO contains $n$
primitive GTOs, whose contraction coefficients and exponents are selected so
that the contraction approximates an STO. STO-$n$G bases are therefore minimal,
single-$\zeta$ basis sets; STO-3G is the most familiar member.
\begin{center}
  \includegraphics[width=0.42\textwidth]{STO-3G}\hspace{1em}
  \includegraphics[width=0.42\textwidth]{STO}
\end{center}

\paragraph{Pople basis sets.}

The notation explicitly records the contractions and augmentations. In
6-31+G*, the core function is contracted from six primitives; the valence
space is double-$\zeta$, with contractions of three and one primitives; the
``+'' adds diffuse functions to non-hydrogen atoms; and ``*'' adds polarization
functions to non-hydrogen atoms, conventionally a $d$ shell. In 6-311++G**,
the valence space is triple-$\zeta$ with contractions $3,1,1$; the second
``+'' adds diffuse functions to hydrogen; and the second ``*'' adds
polarization functions to hydrogen. Finally, 6-311G(2df,2pd) adds two $d$ and
one $f$ polarization shells to non-hydrogen atoms, and two $p$ and one $d$
shell to hydrogen.

\paragraph{Dunning correlation-consistent basis sets.}

The main sequence is
\begin{equation}
  \boxed{\text{cc-pVXZ},\qquad X=\mathrm{D,T,Q,5,6}}.
\end{equation}
These bases were designed for correlated calculations, with balanced layers
of polarization functions and a systematically increasing valence space. A
standard correlation-energy extrapolation is
\begin{equation}
  \Ec(X)\approx\Ec^\text{CBS}+AX^{-3}.
\end{equation}
Diffuse augmentation gives aug-cc-pVXZ, d-aug-cc-pVXZ, and t-aug-cc-pVXZ.
Core correlation is described by cc-pCVXZ, optimized for core--core and
core--valence correlation, or by cc-pwCVXZ, which weights the core--valence
contribution.

For a second-row atom, the valence correlation-consistent sequence is
\begin{center}
\begin{tabular}{c@{\qquad}c}
\toprule
Basis & shell content\\
\midrule
cc-pVDZ & $3s2p1d$\\
cc-pVTZ & $4s3p2d1f$\\
cc-pVQZ & $5s4p3d2f1g$\\
aug-cc-pVDZ & $4s3p2d$\\
aug-cc-pVTZ & $5s4p3d2f$\\
aug-cc-pVQZ & $6s5p4d3f2g$\\
d-aug-cc-pVDZ & $5s4p3d$\\
d-aug-cc-pVTZ & $6s5p4d3f$\\
d-aug-cc-pVQZ & $7s6p5d4f3g$\\
\bottomrule
\end{tabular}
\end{center}
The successive entries make the layered construction explicit. For the
core--valence sequence, the corresponding decomposition is
\begin{align}
 \text{cc-pVDZ: }&3s2p1d=1s+2s2p1d,\\
 \text{cc-pVTZ: }&4s3p2d1f=1s+3s3p2s1f,\\
 \text{cc-pVQZ: }&5s4p3d2f1g=1s+4s4p3d2f1g,\\
 \text{cc-pCVDZ: }&4s3p1d=2s1p+2s2p1d,\\
 \text{cc-pCVTZ: }&6s5p3d1f=3s2p1d+3s3p2d1f,\\
 \text{cc-pCVQZ: }&8s7p5d3f1g=4s3p2d1f+4s4p3d2f1g.
\end{align}

\paragraph{Jensen polarization-consistent basis sets.}

The sequence has shell contents
\begin{equation}
  \text{pc-0}:3s2p,quad
  \text{pc-1}:3s2p1d,quad
  \text{pc-2}:4s3p2d1f,quad
  \text{pc-3}:6s5p4d2f1g.
\end{equation}
Thus pc-0 is unpolarized, pc-1 adds one polarization layer, and the later
members add systematically more radial and angular flexibility. These basis
sets were designed for HF and Kohn--Sham calculations, with higher angular
momenta selected according to their energetic importance. Diffuse variants
are denoted aug-pc-$n$. Their HF convergence is described by
\begin{equation}
  \boxed{E_\text{HF}(L)=E_\text{HF}^\text{CBS}
  +A(L+1)\exp(-B\sqrt L)}.
\end{equation}

Gaussian bases are compact, atom centered, chemically transferable, supported
by analytic integral technology, and available in systematic hierarchies.
Their principal limitations are the absent nuclear cusp, the incorrect
long-range Gaussian decay, basis-set superposition error, and possible linear
dependence when many diffuse functions are used. There is consequently no
single best Gaussian basis: the appropriate choice depends on the property,
the electronic-structure method, and the desired accuracy.

\subsubsection{Basis-set superposition error}

Let $E_X^Y$ denote the energy of system $X$ evaluated in basis $Y$. In a
dimer $AB$, each fragment can borrow the basis functions of the other fragment
as ghost functions. The variational principle therefore gives
\begin{equation}
  E_A^{AB}\leq E_A^A,qquad E_B^{AB}\leq E_B^B,
\end{equation}
which makes the uncorrected interaction energy
\begin{equation}
  \Delta E_\text{int}=E_{AB}^{AB}-E_A^A-E_B^B
\end{equation}
artificially too attractive. The Boys--Bernardi counterpoise correction uses
the full dimer basis for every term,
\begin{equation}
  \boxed{\Delta E_\text{int}^\text{CP}
  =E_{AB}^{AB}-E_A^{AB}-E_B^{AB}},
\end{equation}
and the basis-set superposition error is
\begin{equation}
  \boxed{\mathrm{BSSE}
  =\Delta E_\text{int}^\text{CP}-\Delta E_\text{int}
  =(E_A^A-E_A^{AB})+(E_B^B-E_B^{AB})\geq0}.
\end{equation}

\subsection{Effective core potentials}

The key idea of an effective core potential is to treat only the chemically
relevant electrons explicitly. It replaces the nuclear attraction and the
chemically inactive core electrons by an effective one-electron operator,
\begin{equation}
 \hh(\br)=-\frac{\nabla^2}{2}
 -\sum_A\frac{Z_A}{\abs{\br-\bR_A}}
 \longrightarrow
 \boxed{\hh^\text{ECP}(\br)=-\frac{\nabla^2}{2}
 -\sum_A\frac{\green{Q_A}}{\abs{\br-\bR_A}}
 +\sum_A\hat U_A^\text{ECP}(\br)},
 \qquad \green{Q_A}=Z_A-N_A^\text{core}.
\end{equation}
This lowers the cost by reducing the number of explicit electrons and removing
the basis functions needed for core orbitals. It can incorporate
scalar-relativistic effects and, with a suitable operator, spin--orbit
coupling, and it can be used in both molecular and periodic calculations with
Gaussian or plane-wave bases. The price is a new class of one-electron
integrals that can be difficult to evaluate; core properties are not described
explicitly, and the accuracy depends on the ECP quality.

A common semilocal form is
\begin{equation}
  U^\text{ECP}=U_L(r)+
  \sum_{\ell=0}^{L-1}\sum_{m=-\ell}^{\ell}
  \ket{Y_{\ell m}}U_\ell(r)\bra{Y_{\ell m}}.
\end{equation}
The first term is local, or unprojected; the remaining angular-momentum
channels are nonlocal, or projected. Their radial parts are generally expanded
as
\begin{equation}
  U_\ell(r)=\sum_kD_{\ell k}r^{n_{\ell k}}
  \exp(-\eta_{\ell k}r^2),
  \qquad n_{\ell k}\in\{-2,-1,0\}.
\end{equation}
The $D_{\ell k}$ and $\eta_{\ell k}$ are fitted parameters. Although these are
one-electron operators, their angular projectors make the integrals nontrivial.

For silver, a large-core ECP removes 36 electrons and treats explicitly the
$4d^{10}5s^1$ shells, leaving 11 electrons. A small-core ECP removes 28
electrons and retains $4s^24p^64d^{10}5s^1$, leaving 19 explicit electrons;
the semicore $4s$ and $4p$ shells are therefore retained, providing more
flexibility at higher cost. Common Gaussian-basis ECP families include:
\begin{itemize}
 \item LANL and SBK/SBKJC, established families in which LANL2DZ combines an
 ECP with a double-$\zeta$ valence basis;
 \item \href{https://www.tc.uni-koeln.de/PP/clickpse.en.html}
 {Stuttgart/Cologne energy-consistent pseudopotentials}, offered with different
 core sizes and dedicated bases;
 \item def2-ECPs, paired with the def2 bases of Weigend and Ahlrichs;
 \item Burkatzki--Filippi--Dolg (BFD) potentials, developed for quantum Monte
 Carlo calculations;
 \item \href{https://pseudopotentiallibrary.org}{correlation-consistent ECPs
 (ccECPs)}, developed for high-accuracy correlated many-body methods.
\end{itemize}
Most ECPs have dedicated basis sets.

\subsection{Properties of Gaussian functions}

The central algebraic fact behind Gaussian integral technology is the
\red{Gaussian product rule}: the product of two Gaussians is another Gaussian.
If
\begin{equation}
 G_{\red{\alpha},\red{\bA}}(\br)
 =\exp[-\red{\alpha}\abs{\br-\red{\bA}}^2],
 \qquad
 G_{\blue{\beta},\blue{\bB}}(\br)
 =\exp[-\blue{\beta}\abs{\br-\blue{\bB}}^2],
\end{equation}
then
\begin{equation}
 \boxed{G_{\red{\alpha},\red{\bA}}(\br)
 G_{\blue{\beta},\blue{\bB}}(\br)
 =K_{\violet{\zeta},\violet{\bP}}
 G_{\violet{\zeta},\violet{\bP}}(\br)},
\end{equation}
with
\begin{equation}
 \violet{\zeta}=\red{\alpha}+\blue{\beta},
 \qquad
 \violet{\bP}=\frac{\red{\alpha\bA}+\blue{\beta\bB}}
 {\red{\alpha}+\blue{\beta}},
\end{equation}
and
\begin{equation}
 K_{\violet{\zeta},\violet{\bP}}
 =\exp[-\frac{\red{\alpha}\blue{\beta}}
 {\red{\alpha}+\blue{\beta}}\abs{\red{\bA}-\blue{\bB}}^2].
\end{equation}
The prefactor therefore decays exponentially with the distance between the
two centers.

For two $s$ Gaussians, the overlap illustrates Cartesian factorization:
\begin{align}
 \int G_{\red{\alpha},\red{\bA}}(\br)
 G_{\blue{\beta},\blue{\bB}}(\br)\dd\br
 &=K_{\violet{\zeta},\violet{\bP}}
   \int G_{\violet{\zeta},\violet{\bP}}(\br)\dd\br\\
 &=K_{\violet{\zeta},\violet{\bP}}
 \qty(\int_{-\infty}^{+\infty}e^{-\violet{\zeta}x^2}\dd x)
 \qty(\int_{-\infty}^{+\infty}e^{-\violet{\zeta}y^2}\dd y)
 \qty(\int_{-\infty}^{+\infty}e^{-\violet{\zeta}z^2}\dd z)\\
 &=K_{\violet{\zeta},\violet{\bP}}
 \qty(\frac{\pi}{\violet{\zeta}})^{3/2}.
\end{align}

\subsubsection{One-electron Gaussian integrals}

The Laplacian of an $s$ Gaussian is
\begin{equation}
 -\frac12\laplacian G_{\blue{\beta},\blue{\bB}}(\br)
 =\qty[3\blue{\beta}-2\blue{\beta}^2
 \abs{\br-\blue{\bB}}^2]G_{\blue{\beta},\blue{\bB}}(\br).
\end{equation}
Consequently, its kinetic-energy integral is obtained without omitting the
intermediate step:
\begin{align}
 -\frac12\int G_{\red{\alpha},\red{\bA}}(\br)\laplacian
 G_{\blue{\beta},\blue{\bB}}(\br)\dd\br
 &=\int G_{\red{\alpha},\red{\bA}}(\br)
 \qty[3\blue{\beta}-2\blue{\beta}^2
 \abs{\br-\blue{\bB}}^2]
 G_{\blue{\beta},\blue{\bB}}(\br)\dd\br\\
 &=\frac{\red{\alpha}\blue{\beta}}{\violet{\zeta}}
 \qty[3-2\frac{\red{\alpha}\blue{\beta}}{\violet{\zeta}}
 \abs{\red{\bA}-\blue{\bB}}^2]
 K_{\violet{\zeta},\violet{\bP}}
 \qty(\frac{\pi}{\violet{\zeta}})^{3/2}.
\end{align}

For a nucleus of charge $Z_C$ located at $\bC$, the nuclear-attraction
integral is
\begin{align}
 -Z_C\int G_{\red{\alpha},\red{\bA}}(\br)
 \frac{1}{\abs{\br-\bC}}
 G_{\blue{\beta},\blue{\bB}}(\br)\dd\br
 &=-Z_CK_{\violet{\zeta},\violet{\bP}}
 \int\frac{e^{-\violet{\zeta}\abs{\br-\violet{\bP}}^2}}
 {\abs{\br-\bC}}\dd\br\\
 &=-\frac{2Z_C}{\sqrt\pi}K_{\violet{\zeta},\violet{\bP}}
 \int_0^\infty\dd t\int
 e^{-\violet{\zeta}\abs{\br-\violet{\bP}}^2}
 e^{-t^2\abs{\br-\bC}^2}\dd\br\\
 &=-\frac{2\pi Z_C}{\violet{\zeta}}
 K_{\violet{\zeta},\violet{\bP}}
 F_0\qty(\violet{\zeta}\abs{\violet{\bP}-\bC}^2).
\end{align}

\subsection{Four-center electron-repulsion integrals}

Applying the Gaussian product rule to both electron coordinates gives
\begin{align}
 \rbraket{\red{\ba}\blue{\bb}}{\orange{\bc}\green{\bd}}
 &=\iint G_{\red{\alpha},\red{\bA}}(\br_1)
 G_{\blue{\beta},\blue{\bB}}(\br_1)\frac{1}{r_{12}}
 G_{\orange{\gamma},\orange{\bC}}(\br_2)
 G_{\green{\delta},\green{\bD}}(\br_2)\dd\br_1\dd\br_2\\
 &=K_{\violet{\zeta},\violet{\bP}}
 K_{\purple{\eta},\purple{\bQ}}
 \iint G_{\violet{\zeta},\violet{\bP}}(\br_1)
 \frac{1}{r_{12}}G_{\purple{\eta},\purple{\bQ}}(\br_2)
 \dd\br_1\dd\br_2.
\end{align}
There are formally $\order{K^4}$ ERIs, although only $\order{K^2}$ are
significant in a sufficiently large, spatially extended system.

The notation distinguishes individual functions from shells and primitives
from contractions. Thus $\sbra{\red{\ba}\blue{\bb}}$ is one primitive bra,
whereas $\sbra{\red a\blue b}$ is a complete primitive shell pair;
$\sbraket{\red{\ba}\blue{\bb}}{\orange{\bc}\green{\bd}}$ is one primitive
integral, whereas $\sbraket{\red a\blue b}{\orange c\green d}$ is a primitive
shell quartet. Likewise, $\rbra{\red{\ba}\blue{\bb}}$ and
$\rbra{\red a\blue b}$ denote a contracted bra and contracted shell pair,
and $\rbraket{\red{\ba}\blue{\bb}}{\orange{\bc}\green{\bd}}$ and
$\rbraket{\red a\blue b}{\orange c\green d}$ denote a contracted integral
and a contracted shell quartet.

\subsubsection{The fundamental $ssss$ integral}

For four primitive $s$ functions,
\begin{align}
 \sbraket{\bs\bs}{\bs\bs}
 &=K_{\violet{\zeta},\violet{\bP}}
 K_{\purple{\eta},\purple{\bQ}}
 \iint G_{\violet{\zeta},\violet{\bP}}(\br_1)
 \frac{1}{r_{12}}G_{\purple{\eta},\purple{\bQ}}(\br_2)
 \dd\br_1\dd\br_2\\
 &=K_{\violet{\zeta},\violet{\bP}}
 K_{\purple{\eta},\purple{\bQ}}
 \iint e^{-\violet{\zeta}\abs{\br_1-\violet{\bP}}^2}
 \frac{1}{r_{12}}
 e^{-\purple{\eta}\abs{\br_2-\purple{\bQ}}^2}
 \dd\br_1\dd\br_2.
\end{align}
The Laplace transform
\begin{equation}
 \frac{1}{r_{12}}=\frac{2}{\sqrt\pi}\int_0^\infty
 e^{-t^2r_{12}^2}\dd t
 =\frac{2}{\sqrt\pi}\int_0^\infty
 e^{-t^2(x_1-x_2)^2}e^{-t^2(y_1-y_2)^2}
 e^{-t^2(z_1-z_2)^2}\dd t
\end{equation}
turns the six-dimensional integral into three Cartesian factors:
\begin{align}
 \sbraket{\bs\bs}{\bs\bs}
 ={}&\frac{2}{\sqrt\pi}K_{\violet{\zeta},\violet{\bP}}
 K_{\purple{\eta},\purple{\bQ}}\int_0^\infty\dd t
 \qty[\iint e^{-\violet{\zeta}(x_1-\violet P_x)^2}
 e^{-t^2(x_1-x_2)^2}e^{-\purple{\eta}(x_2-\purple Q_x)^2}
 \dd x_1\dd x_2]\\
 &\cdot\qty[\iint e^{-\violet{\zeta}(y_1-\violet P_y)^2}
 e^{-t^2(y_1-y_2)^2}e^{-\purple{\eta}(y_2-\purple Q_y)^2}
 \dd y_1\dd y_2]\\
 &\cdot\qty[\iint e^{-\violet{\zeta}(z_1-\violet P_z)^2}
 e^{-t^2(z_1-z_2)^2}e^{-\purple{\eta}(z_2-\purple Q_z)^2}
 \dd z_1\dd z_2].
\end{align}

Defining
\begin{equation}
 \vartheta=\frac{\violet{\zeta}\purple{\eta}}
 {\violet{\zeta}+\purple{\eta}},
\end{equation}
the Cartesian integrations yield, step by step,
\begin{align}
 \sbraket{\bs\bs}{\bs\bs}
 ={}&\frac{2}{\sqrt\pi}K_{\violet{\zeta},\violet{\bP}}
 K_{\purple{\eta},\purple{\bQ}}\int_0^\infty
 \prod_{i=x,y,z}\qty[
 \frac{\pi}{\sqrt{\violet{\zeta}\purple{\eta}}}
 \frac{\exp[-\vartheta(\violet P_i-\purple Q_i)^2
 t^2/(t^2+\vartheta)]}
 {\sqrt{1+t^2/\vartheta}}]\dd t\\
 ={}&\frac{2\pi^{5/2}}
 {(\violet{\zeta}\purple{\eta})^{3/2}}
 K_{\violet{\zeta},\violet{\bP}}K_{\purple{\eta},\purple{\bQ}}
 \int_0^\infty
 \frac{\exp[-\vartheta\abs{\violet{\bP}-\purple{\bQ}}^2
 t^2/(t^2+\vartheta)]}{(1+t^2/\vartheta)^{3/2}}\dd t.
\end{align}
With the substitution $t^2=\vartheta u^2/(1-u^2)$,
\begin{align}
 \sbraket{\bs\bs}{\bs\bs}
 ={}&\sqrt{\frac{4\vartheta}{\pi}}
 \qty(\frac{\pi}{\violet{\zeta}})^{3/2}
 K_{\violet{\zeta},\violet{\bP}}
 \qty(\frac{\pi}{\purple{\eta}})^{3/2}
 K_{\purple{\eta},\purple{\bQ}}
 \underbrace{\int_0^1e^{-\vartheta
 \abs{\violet{\bP}-\purple{\bQ}}^2u^2}\dd u}
 _{F_0(\vartheta\abs{\violet{\bP}-\purple{\bQ}}^2)}.
\end{align}
The fundamental auxiliary notation is therefore
\begin{equation}
 \boxed{
 \sbraket{\bO\bO}{\bO\bO}^{(0)}
 =\sqrt{\frac{4\vartheta}{\pi}}
 \qty(\frac{\pi}{\violet{\zeta}})^{3/2}K_{\violet{\zeta},\violet{\bP}}
 \qty(\frac{\pi}{\purple{\eta}})^{3/2}K_{\purple{\eta},\purple{\bQ}}
 F_0\qty(\vartheta\abs{\violet{\bP}-\purple{\bQ}}^2)}.
\end{equation}

\subsubsection{Boys functions and auxiliary integrals}

The generalized Boys function is
\begin{equation}
 F_m(x)=\int_0^1t^{2m}\exp(-xt^2)\dd t,
 \qquad
 F_0(x)=\sqrt{\frac{\pi}{4x}}\erf(\sqrt x).
\end{equation}
In implementations, $F_m(x)$ is evaluated from pretabulated grid values and
a Taylor expansion, often with a Chebyshev interpolation table. Lower orders
are then obtained by the downward recurrence
\begin{equation}
 F_m(x)=\frac{2xF_{m+1}(x)+\exp(-x)}{2m+1}.
\end{equation}
This downward recurrence is numerically stable, whereas upward recursion can
become unstable for small $x$. At large $x$,
\begin{equation}
 F_m(x)\sim\frac{(2m-1)!!}{2^{m+1}}
 \sqrt{\frac{\pi}{x^{2m+1}}}.
\end{equation}

The physical $ssss$ integral contains only $F_0$, but differentiation to
generate angular momentum produces higher Boys orders. It is therefore useful
to define the complete hierarchy
\begin{equation}
 \boxed{
 \sbraket{\bO\bO}{\bO\bO}^{(\red m)}
 =\sqrt{\frac{4\vartheta}{\pi}}
 \qty(\frac{\pi}{\violet{\zeta}})^{3/2}K_{\violet{\zeta},\violet{\bP}}
 \qty(\frac{\pi}{\purple{\eta}})^{3/2}K_{\purple{\eta},\purple{\bQ}}
 F_{\red m}\qty(\vartheta\abs{\violet{\bP}-\purple{\bQ}}^2)}.
\end{equation}
The $m=0$ member is the physical $ssss$ ERI; the $m>0$ members are auxiliary
intermediates. Angular-momentum recurrences couple orders $m$ and $m+1$.
For total angular momentum $L=a+b+c+d$, fundamental auxiliary integrals are
needed through $m=L$.

\subsubsection{Obara--Saika recurrences}

For a Cartesian Gaussian, the key identity is
\begin{equation}
 \pdv{\sket{\red{\ba}}}{A_i}
 =2\alpha\sket{\red{\ba}+\bI_i}
 -a_i\sket{\red{\ba}-\bI_i},
 \qquad
 \boxed{\sket{\red{\ba}+\bI_i}
 =\frac{1}{2\alpha}\qty[
 \pdv{\sket{\red{\ba}}}{A_i}
 +a_i\sket{\red{\ba}-\bI_i}]}.
\end{equation}
Applying this identity to auxiliary ERIs creates the Obara--Saika vertical
recurrence relations (VRRs), so that integrals are built recursively rather
than independently:
\begin{equation}
 \boxed{\sbraket{00}{00}^{(m)}\longrightarrow
 \sbraket{\red a0}{00}^{(m)}\longrightarrow
 \sbraket{\red a0}{\orange c0}^{(m)}}.
\end{equation}

On the bra pair, angular momentum is built on $\bA$ according to
\begin{align}
 \sbraket{(\red{\ba}+\bI_i)\bO}{\bO\bO}^{(m)}
 ={}&(P_i-A_i)\sbraket{\red{\ba}\bO}{\bO\bO}^{(m)}
 +\frac{\eta}{\zeta+\eta}(P_i-Q_i)
 \sbraket{\red{\ba}\bO}{\bO\bO}^{(m+1)}\notag\\
 &+\frac{a_i}{2\zeta}\qty[
 \sbraket{(\red{\ba}-\bI_i)\bO}{\bO\bO}^{(m)}
 -\frac{\eta}{\zeta+\eta}
 \sbraket{(\red{\ba}-\bI_i)\bO}{\bO\bO}^{(m+1)}].
\end{align}
On the ket pair, angular momentum is built on $\bC$ according to
\begin{align}
 \sbraket{\red{\ba}\bO}{(\orange{\bc}+\bI_i)\bO}^{(m)}
 ={}&(Q_i-C_i)\sbraket{\red{\ba}\bO}{\orange{\bc}\bO}^{(m)}
 +\frac{\zeta}{\zeta+\eta}(Q_i-P_i)
 \sbraket{\red{\ba}\bO}{\orange{\bc}\bO}^{(m+1)}\notag\\
 &+\frac{c_i}{2\eta}\qty[
 \sbraket{\red{\ba}\bO}{(\orange{\bc}-\bI_i)\bO}^{(m)}
 -\frac{\zeta}{\zeta+\eta}
 \sbraket{\red{\ba}\bO}{(\orange{\bc}-\bI_i)\bO}^{(m+1)}]\notag\\
 &+\frac{a_i}{2(\zeta+\eta)}
 \sbraket{(\red{\ba}-\bI_i)\bO}{\orange{\bc}\bO}^{(m+1)}.
\end{align}

The intended construction is
\begin{equation}
 \boxed{
 \sbraket{00}{00}^{(m)}
 \xrightarrow{\red{\mathrm{VRR}}}
 \sbraket{\red a0}{\orange c0}^{(m)}
 \xrightarrow{\orange{\mathrm{contraction}}}
 \rbraket{\red a0}{\orange c0}
 \xrightarrow{\blue{\mathrm{HRR}}}
 \rbraket{\red a\blue b}{\orange c\green d}}.
\end{equation}
The VRRs create angular momentum on one center of each Gaussian pair;
contraction combines primitive integrals; and the horizontal recurrence
relations (HRRs) redistribute angular momentum within each pair.

Once angular momentum has been generated on one center, it is shifted to the
other through
\begin{equation}
 \boxed{
 \sbraket{\red{\ba}(\blue{\bb}+\bI_i)}
 {\orange{\bc}\green{\bd}}
 =\sbraket{(\red{\ba}+\bI_i)\blue{\bb}}
 {\orange{\bc}\green{\bd}}
 +(A_i-B_i)\sbraket{\red{\ba}\blue{\bb}}
 {\orange{\bc}\green{\bd}}}.
\end{equation}
An analogous relation transfers angular momentum from $\bC$ to $\bD$.
Because the HRR depends only on $A_i-B_i$, it is independent of the Gaussian
exponents and applies directly to contracted integrals:
\begin{equation}
 \boxed{
 \rbraket{\red{\ba}(\blue{\bb}+\bI_i)}
 {\orange{\bc}\green{\bd}}
 =\rbraket{(\red{\ba}+\bI_i)\blue{\bb}}
 {\orange{\bc}\green{\bd}}
 +(A_i-B_i)\rbraket{\red{\ba}\blue{\bb}}
 {\orange{\bc}\green{\bd}}}.
\end{equation}

The resulting late-contraction algorithm retains every stage:
\begin{enumerate}
 \item compute the required Boys functions $F_m(T)$;
 \item form the fundamental auxiliary integrals
 $\sbraket{00}{00}^{(m)}$;
 \item use the VRRs to build angular momentum on $\bA$ and $\bC$;
 \item contract the primitive integrals;
 \item use exponent-independent HRRs to transfer angular momentum
 $\bA\rightarrow\bB$ and $\bC\rightarrow\bD$.
\end{enumerate}
Thus the expensive exponent-dependent recurrences are performed before
contraction, whereas the cheap exponent-independent HRRs are performed after
contraction.

\subsection{Bounds and screening}

A useful ERI upper bound must be simple and inexpensive, strong enough to be a
good estimate, and scaling consistent, meaning
$N_\text{sig}=\order{N_\text{upper}}$. The Cauchy--Schwarz inequality gives
\begin{equation}
 \abs{\rbraket{\red{\ba}\blue{\bb}}{\orange{\bc}\green{\bd}}}
 \leq
 \sqrt{\rbraket{\red{\ba}\blue{\bb}}{\red{\ba}\blue{\bb}}}
 \sqrt{\rbraket{\orange{\bc}\green{\bd}}{\orange{\bc}\green{\bd}}},
\end{equation}
or, in terms of the Gaussian product centers,
\begin{equation}
 \abs{\rbraket{\violet{\bP}}{\purple{\bQ}}}
 \leq\sqrt{\rbraket{\violet{\bP}}{\violet{\bP}}}
 \sqrt{\rbraket{\purple{\bQ}}{\purple{\bQ}}}.
\end{equation}
The integral $\rbraket{\violet{\bP}}{\purple{\bQ}}$ should decay as
$1/\abs{\violet{\bP}-\purple{\bQ}}$ or faster. The product on the right has
no information about the separation between the two charge distributions, so
the Cauchy--Schwarz bound is rigorous and cheap but not tight for well-separated
shell pairs. Even so, a cheap upper bound avoids evaluating negligible ERIs.

At the individual-integral level, only $\order{K^2}$ diagonal factors
$\rbraket{\red{\ba}\blue{\bb}}{\red{\ba}\blue{\bb}}$ are required to bound
$\order{K^4}$ four-center integrals. More efficiently, define a shell-pair
bound
\begin{equation}
 Q_{\red a\blue b}
 =\max_{\substack{\red{\ba}\in a\\\blue{\bb}\in b}}
 \sqrt{\rbraket{\red{\ba}\blue{\bb}}{\red{\ba}\blue{\bb}}}.
\end{equation}
Then an entire shell class obeys
\begin{equation}
 \abs{\rbraket{\red a\blue b}{\orange c\green d}}
 \leq Q_{\red a\blue b}Q_{\orange c\green d}.
\end{equation}
If
\begin{equation}
 Q_{\red a\blue b}Q_{\orange c\green d}<\tau,
\end{equation}
the entire shell quartet can be discarded. This is consequential: even the
simple $\rbraket{dd}{dd}$ class contains 1296 Cartesian integrals.

The number of significant integrals may be represented by
\begin{equation}
 \rbraket{\red{\ba}\blue{\bb}}{\orange{\bc}\green{\bd}}
 \equiv\rmel{\red{\ba}\blue{\bb}}{\chO_2}
 {\orange{\bc}\green{\bd}},
 \qquad N_\text{sig}=cN^\alpha.
\end{equation}
The numerical examples in the slides contrast the long-range Coulomb operator
with a short-range Gaussian operator:
\begin{center}
\small
\begin{tabular}{lrrrrr}
\toprule
& &\multicolumn{2}{c}{\red{$\chO=r_{12}^{-1}$}}
&\multicolumn{2}{c}{\orange{$\chO=e^{-r_{12}^2}$}}\\
\cmidrule(lr){3-4}\cmidrule(lr){5-6}
Molecule&$N$&$N_\text{sig}$&$\alpha$&$N_\text{sig}$&$\alpha$\\
\midrule
propene&12&1,625&---&1,650&---\\
butadiene&16&5,020&3.9&5,020&3.9\\
hexatriene&24&24,034&3.9&23,670&3.8\\
octatetraene&32&63,818&3.4&52,808&2.8\\
decapentaene&40&119,948&2.8&81,404&1.9\\
dodecaexaene&48&192,059&2.6&109,965&1.6\\
\bottomrule
\end{tabular}
\end{center}
The contrast shows the more rapid approach to low effective scaling for a
short-range operator.

The complete recipe has three ingredients. The ``cake'' is the set of target
two-electron integrals $\rbraket{\ba\bb}{\bc\bd}$. Ingredient 1 is the set of
fundamental integrals $\rbraket{\bO\bO}{\bO\bO}^{(m)}$; ingredient 2 is the
set of recurrence relations, schematically
$\rbraket{(\ba+\bI_i)\bO}{\bO\bO}
=\rbraket{\ba\bO}{\bO\bO}
+\rbraket{(\ba-\bI_i)\bO}{\bO\bO}$; and ingredient 3 is an upper bound
$\abs{\rbraket{\ba\bb}{\bc\bd}}\leq B$.

\subsubsection{Direct SCF}

The Coulomb and exchange matrices are
\begin{align}
 J_{\mu\nu}&=\sum_{\lambda\sigma}P_{\lambda\sigma}
 \rbraket{\mu\nu}{\lambda\sigma},
 &
 K_{\mu\nu}&=\sum_{\lambda\sigma}P_{\lambda\sigma}
 \rbraket{\mu\lambda}{\nu\sigma}.
\end{align}
Direct SCF trades disk storage for repeated integral evaluation. Instead of
storing all $\order{K^4}$ ERIs, it performs the following cycle:
\begin{enumerate}
 \item loop over shell quartets $\rbraket{\mu\nu}{\lambda\sigma}$;
 \item use an upper bound to discard negligible quartets;
 \item compute the significant ERIs;
 \item immediately accumulate their contributions to $\bJ$ and $\bK$;
 \item discard the ERIs.
\end{enumerate}
The density-matrix elements can themselves be included in the screening
criterion. In compact form, the algorithm is
\begin{equation}
 \boxed{\text{screen}\longrightarrow\text{compute}\longrightarrow
 \text{accumulate}\longrightarrow\text{discard}}.
\end{equation}

\subsection{Resolution of the identity}

An ERI is the Coulomb interaction between two pair densities,
\begin{equation}
 \rbraket{\mu\nu}{\lambda\sigma}
 =\iint
 \underbrace{\phi_\mu(\br_1)\phi_\nu(\br_1)}_{\rho_{\mu\nu}(\br_1)}
 \frac{1}{r_{12}}
 \underbrace{\phi_\lambda(\br_2)\phi_\sigma(\br_2)}
 _{\rho_{\lambda\sigma}(\br_2)}\dd\br_1\dd\br_2.
\end{equation}
With $K$ orbital basis functions, there are formally $\order{K^2}$ different
pair densities. Resolution of the identity, or density fitting, replaces this
large product space by a compact auxiliary basis
$\{\varphi_A^\text{aux}\}$:
\begin{equation}
 \boxed{\rho_{\mu\nu}(\br)\approx
 \widetilde\rho_{\mu\nu}(\br)
 =\sum_AC_{\mu\nu}^A\varphi_A^\text{aux}(\br)},
 \qquad K_\text{aux}=\order{K}.
\end{equation}

\subsubsection{The Coulomb metric}

Define the Coulomb inner product
\begin{equation}
 \rbraket{f}{g}=\iint f(\br_1)\frac{1}{r_{12}}g(\br_2)
 \dd\br_1\dd\br_2.
\end{equation}
The optimal fitting coefficients minimize the Coulomb norm of the error,
\begin{equation}
 \Delta_{\mu\nu}=\rbraket{
ho_{\mu\nu}-\widetilde\rho_{\mu\nu}}
 {\rho_{\mu\nu}-\widetilde\rho_{\mu\nu}}.
\end{equation}
Differentiation with respect to $C_{\mu\nu}^A$ yields the least-squares
equations
\begin{equation}
 \sum_B\underbrace{\rbraket{A}{B}}_{V_{AB}}C_{\mu\nu}^B
 =\underbrace{\rbraket{A}{\mu\nu}}_{\text{three-center integral}},
 \qquad
 \boxed{C_{\mu\nu}^A=\sum_B(\bV^{-1})_{AB}\rbraket{B}{\mu\nu}}.
\end{equation}
Thus the Coulomb metric determines the optimal auxiliary expansion.

Inserting the fitted pair densities converts the four-center ERI to
three-center quantities:
\begin{equation}
 \boxed{\rbraket{\mu\nu}{\lambda\sigma}
 \approx\sum_{AB}\rbraket{\mu\nu}{A}(\bV^{-1})_{AB}
 \rbraket{B}{\lambda\sigma}}.
\end{equation}
With the symmetrically orthonormalized quantities
\begin{equation}
 B_{\mu\nu}^P=\sum_A\rbraket{\mu\nu}{A}(\bV^{-1/2})_{AP},
\end{equation}
the factorization becomes
\begin{equation}
 \boxed{\rbraket{\mu\nu}{\lambda\sigma}
 \approx\sum_PB_{\mu\nu}^PB_{\lambda\sigma}^P}.
\end{equation}
The four-index tensor has thereby been replaced by products of three-index
tensors.

\subsubsection{Low-rank representations and RI in practice}

In RI/DF, the AO products are expanded in a predefined auxiliary basis,
\begin{equation}
 \chi_\mu(\br)\chi_\nu(\br)
 \approx\sum_AC_{\mu\nu}^A\chi_A^\text{aux}(\br)
 \quad\Longrightarrow\quad
 \boxed{\rbraket{\mu\nu}{\lambda\sigma}
 \approx\sum_PB_{\mu\nu}^PB_{\lambda\sigma}^P}.
\end{equation}
The quality of the approximation therefore depends on the auxiliary basis.
Alternatively, the ERI tensor
$\bV_{(\mu\nu),(\lambda\sigma)}=\rbraket{\mu\nu}{\lambda\sigma}$ is a
positive-semidefinite matrix in the compound indices. A pivoted Cholesky
decomposition constructs an auxiliary space adaptively,
\begin{equation}
 \boxed{\rbraket{\mu\nu}{\lambda\sigma}
 \approx\sum_{P=1}^{N_\text{CD}}L_{\mu\nu}^PL_{\lambda\sigma}^P},
\end{equation}
until the residual is below a prescribed threshold $\tau_\text{CD}$. RI/DF
therefore uses a predefined auxiliary space, whereas Cholesky decomposition
constructs one adaptively.

For the Coulomb matrix,
\begin{align}
 J_{\mu\nu}
 &=\sum_{\lambda\sigma}P_{\lambda\sigma}
 \rbraket{\mu\nu}{\lambda\sigma}\\
 &\approx\sum_{AB}\rbraket{\mu\nu}{A}(\bV^{-1})_{AB}
 \underbrace{\sum_{\lambda\sigma}P_{\lambda\sigma}
 \rbraket{B}{\lambda\sigma}}_{d_B}.
\end{align}
The practical sequence is
\begin{equation}
 \boxed{\bP\longrightarrow\{d_A\}\longrightarrow
 \{c_A\}=\bV^{-1}\cdot\{d_A\}\longrightarrow\bJ}.
\end{equation}
The full $\order{K^4}$ ERI tensor need not be stored; only two-center
$\rbraket{A}{B}$ and three-center $\rbraket{\mu\nu}{A}$ integrals are needed;
and optimized auxiliary bases make the fitting error very small. RI is thus
also commonly called density fitting (DF), and it greatly reduces the cost and
memory requirements of Coulomb contractions.
